\documentclass[12pt,a4paper]{article}

\usepackage[utf8]{inputenc}
\usepackage[T1]{fontenc}
\usepackage[french]{babel}
\usepackage{geometry}
\geometry{top=2.5cm, bottom=2.5cm, left=2.5cm, right=2.5cm}

\usepackage{graphicx}
\usepackage{float}
\usepackage{xcolor}
\usepackage{listings}
\usepackage{hyperref}
\usepackage{fancyhdr}
\usepackage{titlesec}
\usepackage{amsmath,amssymb}

% Configuration des métadonnées PDF et liens
\hypersetup{
    colorlinks=true,
    linkcolor=blue!80!black,
    urlcolor=blue!70!black,
    citecolor=blue!80!black
}

% Couleurs pour la coloration syntaxique
\definecolor{codegreen}{rgb}{0,0.6,0}
\definecolor{codegray}{rgb}{0.5,0.5,0.5}
\definecolor{codepurple}{rgb}{0.58,0,0.82}
\definecolor{backcolour}{rgb}{0.97,0.97,0.98}

% Configuration de listings pour le code Python
\lstdefinestyle{pythonstyle}{
    backgroundcolor=\color{backcolour},   
    commentstyle=\color{codegreen}\itshape,
    keywordstyle=\color{blue}\bfseries,
    numberstyle=\tiny\color{codegray},
    stringstyle=\color{codepurple},
    basicstyle=\ttfamily\footnotesize,
    breakatwhitespace=false,         
    breaklines=true,                 
    captionpos=b,                    
    keepspaces=true,                 
    numbers=left,                    
    numbersep=7pt,                  
    showspaces=false,                
    showstringspaces=false,
    showtabs=false,                  
    tabsize=4,
    frame=single,
    rulecolor=\color{gray!30},
    literate=
      {á}{{\'a}}1 {é}{{\'e}}1 {í}{{\'i}}1 {ó}{{\'o}}1 {ú}{{\'u}}1
      {Á}{{\'A}}1 {É}{{\'E}}1 {Í}{{\'I}}1 {Ó}{{\'O}}1 {Ú}{{\'U}}1
      {à}{{\`a}}1 {è}{{\`e}}1 {ì}{{\`i}}1 {ò}{{\`o}}1 {ù}{{\`u}}1
      {À}{{\`A}}1 {È}{{\'E}}1 {Ì}{{\`I}}1 {Ò}{{\`O}}1 {Ù}{{\`U}}1
      {ä}{{\"a}}1 {ë}{{\"e}}1 {ï}{{\"i}}1 {ö}{{\"o}}1 {ü}{{\"u}}1
      {Ä}{{\"A}}1 {Ë}{{\"E}}1 {Ï}{{\"I}}1 {Ö}{{\"O}}1 {Ü}{{\"U}}1
      {â}{{\^a}}1 {ê}{{\^e}}1 {î}{{\^i}}1 {ô}{{\^o}}1 {û}{{\^u}}1
      {Â}{{\^A}}1 {Ê}{{\^E}}1 {Î}{{\^I}}1 {Ô}{{\^O}}1 {Û}{{\^U}}1
      {œ}{{\oe}}1 {Œ}{{\OE}}1 {æ}{{\ae}}1 {Æ}{{\AE}}1 {ç}{{\c{c}}}1
      {Ç}{{\c{C}}}1 {«}{{\guillemotleft}}1 {»}{{\guillemotright}}1
      {°}{{\ensuremath{^\circ}}}1
}

\lstset{style=pythonstyle}

% Configuration des en-têtes et pieds de page
\pagestyle{fancy}
\fancyhf{}
\fancyhead[L]{\small\nouppercase{\leftmark}}
\fancyhead[R]{\small Classification de fleurs -- MobileNetV2}
\fancyfoot[C]{\thepage}
\renewcommand{\headrulewidth}{0.4pt}
\renewcommand{\footrulewidth}{0.4pt}

\title{
    \vspace{-1.5cm}
    \textbf{\Huge Rapport de Projet} \\
    \vspace{0.4cm}
    \textbf{\Large Classification d'Images de Fleurs avec Transfer Learning (MobileNetV2) \& Déploiement Streamlit}
}
\author{\textbf{Auteur : Houdaifa}}
\date{\today}

\begin{document}

\maketitle

\begin{abstract}
Ce rapport présente la conception, l'expérimentation et le déploiement d'une solution complète d'apprentissage profond pour la classification automatique d'images florales issue du jeu de données \textbf{TF Flowers} (3\,670 images réparties en 5 espèces : pissenlit, marguerite, tulipe, tournesol et rose). Le document détaille l'ensemble des cellules du notebook de développement (\texttt{model\_ft\_flowers\_MobileNet.ipynb}), l'intégration d'un réseau convolutif pré-entraîné (\textbf{MobileNetV2}) exploitant le \textit{Transfer Learning}, ainsi que l'application Web interactive développée avec \textbf{Streamlit} (\texttt{app.py}) permettant l'inférence en temps réel via téléversement ou caméra.
\end{abstract}

\vspace{0.5cm}
\tableofcontents
\newpage

\section{Introduction et Contexte du Projet}
La vision par ordinateur appliquée à la botanique présente un défi important en raison des variations d'échelle, d'orientation, d'éclairage et des fortes ressemblances intra-familles. Le projet met en œuvre une démarche méthodique en deux volets principaux :
\begin{enumerate}
    \item \textbf{Phase de modélisation et recherche :} exploration quantitative et qualitative du dataset, mise en place d'un pipeline de prétraitement et d'augmentation de données, benchmark d'une architecture CNN de base, puis mise en œuvre d'un modèle optimisé fondé sur \textbf{MobileNetV2} avec transfert d'apprentissage.
    \item \textbf{Phase de mise en production :} déploiement du réseau au sein d'une interface web ergonomique propulsée par le framework \textbf{Streamlit}, offrant une classification interactive accompagnée d'une estimation de l'incertitude (seuil de confiance) et d'une visualisation des distributions de probabilités.
\end{enumerate}

\newpage
\section{Développement et Analyse du Notebook (\texttt{model\_ft\_flowers\_MobileNet.ipynb})}
Cette section retrace l'intégralité du cycle de travail réalisé dans l'environnement de développement Jupyter / Colab. Chaque sous-section correspond à une cellule du notebook avec son code source, son objectif et son résultat d'exécution.

\subsection{Cellule 0 : Installation des Dépendances}
La première cellule assure l'installation silencieuse et la mise à jour des modules essentiels, en particulier \texttt{tensorflow-datasets} pour l'acquisition et la structuration des jeux de données, ainsi que \texttt{importlib-resources}.

\begin{lstlisting}[language=Python]
!pip install -q tensorflow-datasets
!pip install -U importlib-resources tensorflow-datasets
\end{lstlisting}

\begin{figure}[H]
  \centering
  \includegraphics[width=0.82\textwidth]{1.png}
\end{figure}

\subsection{Cellule 1 : Importation des Bibliothèques et Environnement}
Importation de TensorFlow, TFDS, Matplotlib pour les tracés graphiques et NumPy pour les opérations matricielles. La version de TensorFlow chargée est affichée.

\begin{lstlisting}[language=Python]
import tensorflow as tf
import tensorflow_datasets as tfds
import matplotlib.pyplot as plt
import numpy as np

print("Version TensorFlow :", tf.__version__)
\end{lstlisting}

\begin{figure}[H]
  \centering
  \includegraphics[width=0.82\textwidth]{2.png}
\end{figure}

\subsection{Cellule 2 : Téléchargement et Découpage du Dataset TF Flowers}
Chargement du jeu de données \texttt{tf\_flowers} et division partitionnée des 3\,670 images selon une répartition standard : 80\% pour l'entraînement, 10\% pour la validation et 10\% pour le test. Extraction des noms des 5 classes.

\begin{lstlisting}[language=Python]
(train_ds, val_ds, test_ds), metadata = tfds.load(
    "tf_flowers",
    split=[
        "train[:80%]",
        "train[80%:90%]",
        "train[90%:]"
    ],
    with_info=True,
    as_supervised=True
)

class_names = metadata.features["label"].names

print("Dataset téléchargé avec succès !")
print("Nombre de classes :", metadata.features["label"].num_classes)
print("Nombre d'images :", metadata.splits["train"].num_examples)
\end{lstlisting}

\begin{figure}[H]
  \centering
  \includegraphics[width=0.82\textwidth]{3.png}
\end{figure}

\subsection{Cellule 3 : Analyse Quantitative Exploratoire (EDA)}
Dénombrement des occurrences de chaque classe sur l'ensemble complet des données brutes au moyen d'un compteur \texttt{Counter}. Cette étape permet de vérifier l'équilibre des classes.

\begin{lstlisting}[language=Python]
from collections import Counter

# Charger toutes les images originales pour l'analyse
# On utilise un dataset séparé afin de ne pas modifier
# les données utilisées pour l'entraînement.
ds_all = tfds.load(
    "tf_flowers",
    split="train",
    as_supervised=True
)

# Compter le nombre total d'images et les images par classe
class_counts = Counter()

for image, label in tfds.as_numpy(ds_all):
    class_counts[int(label)] += 1

total_images = sum(class_counts.values())
num_classes = len(class_names)

print("===== ANALYSE QUANTITATIVE =====")
print("Nombre total d'images :", total_images)
print("Nombre de classes :", num_classes)

print("\nNombre d'images par classe :")
for i, name in enumerate(class_names):
    print(f"{name:15s} : {class_counts[i]} images")
\end{lstlisting}

\begin{figure}[H]
  \centering
  \includegraphics[width=0.82\textwidth]{4.png}
\end{figure}

\subsection{Cellule 4 : Visualisation Graphique de la Répartition des Classes}
Génération d'un diagramme en bâtons avec \texttt{matplotlib.pyplot} affichant le nombre exact d'échantillons par catégorie de fleur (\textit{dandelion}, \textit{daisy}, \textit{tulips}, \textit{sunflowers}, \textit{roses}).

\begin{lstlisting}[language=Python]
# Préparer les données pour le graphique
counts = [class_counts[i] for i in range(num_classes)]

plt.figure(figsize=(10, 5))

bars = plt.bar(class_names, counts)

plt.title("Répartition des images par classe - TF Flowers")
plt.xlabel("Catégories de fleurs")
plt.ylabel("Nombre d'images")
plt.xticks(rotation=25)

# Afficher le nombre d'images au-dessus de chaque barre
for bar in bars:
    height = bar.get_height()
    plt.text(
        bar.get_x() + bar.get_width() / 2,
        height,
        str(int(height)),
        ha="center",
        va="bottom"
    )

plt.tight_layout()
plt.show()
\end{lstlisting}

\begin{figure}[H]
  \centering
  \includegraphics[width=0.82\textwidth]{5.png}
\end{figure}

\subsection{Cellule 5 : Analyse Morphologique des Images (Dimensions et Canaux)}
Inspection statistique des résolutions spatiales (largeurs et hauteurs minimales, maximales et les plus fréquentes), des types de codage (\texttt{uint8}) et du nombre de canaux (3 canaux couleur RGB).

\begin{lstlisting}[language=Python]
# Analyser les dimensions et les propriétés des images
image_shapes = Counter()
image_dtypes = Counter()
image_channels = Counter()

min_height, max_height = float("inf"), 0
min_width, max_width = float("inf"), 0

for image, label in tfds.as_numpy(ds_all):
    height, width, channels = image.shape

    # Compter les dimensions rencontrées
    image_shapes[(height, width)] += 1
    image_dtypes[str(image.dtype)] += 1
    image_channels[channels] += 1

    # Rechercher les dimensions minimales et maximales
    min_height = min(min_height, height)
    max_height = max(max_height, height)
    min_width = min(min_width, width)
    max_width = max(max_width, width)

print("===== DIMENSIONS ET FORMAT =====")
print("Dimensions les plus fréquentes :")
for shape, count in image_shapes.most_common(10):
    print(f"{shape[1]} x {shape[0]} : {count} images")

print("\nHauteur minimale :", min_height)
print("Hauteur maximale :", max_height)
print("Largeur minimale :", min_width)
print("Largeur maximale :", max_width)

print("\nTypes de données :", image_dtypes)
print("Nombre de canaux :", image_channels)
\end{lstlisting}

\begin{figure}[H]
  \centering
  \includegraphics[width=0.82\textwidth]{6.png}
\end{figure}

\subsection{Cellule 6 : Prétraitement et Pipeline de Données \texttt{tf.data}}
Homogénéisation de toutes les images à une résolution fixe de $180 \times 180$ pixels, conversion des valeurs de pixels en flottants dans l'intervalle $[0, 1]$, et mise en place des optimisations d'exécution (\texttt{shuffle}, \texttt{batch} de taille 32, et \texttt{prefetch} avec autotuning).

\begin{lstlisting}[language=Python]
# Paramètres du prétraitement
IMG_SIZE = (180, 180)
BATCH_SIZE = 32
NUM_CLASSES = 5

# Fonction de prétraitement des images
def preprocess(image, label):
    # Redimensionner toutes les images à la même résolution
    image = tf.image.resize(image, IMG_SIZE)

    # Convertir les pixels de [0, 255] vers [0, 1]
    image = tf.cast(image, tf.float32) / 255.0

    # Conserver l'étiquette entière de la classe
    return image, label

# Appliquer le prétraitement aux trois ensembles
train_ds = (
    train_ds
    .map(preprocess, num_parallel_calls=tf.data.AUTOTUNE)
    .shuffle(1000)
    .batch(BATCH_SIZE)
    .prefetch(tf.data.AUTOTUNE)
)

val_ds = (
    val_ds
    .map(preprocess, num_parallel_calls=tf.data.AUTOTUNE)
    .batch(BATCH_SIZE)
    .prefetch(tf.data.AUTOTUNE)
)

test_ds = (
    test_ds
    .map(preprocess, num_parallel_calls=tf.data.AUTOTUNE)
    .batch(BATCH_SIZE)
    .prefetch(tf.data.AUTOTUNE)
)

print("Prétraitement terminé.")
\end{lstlisting}

\begin{figure}[H]
  \centering
  \includegraphics[width=0.82\textwidth]{7.png}
\end{figure}

\subsection{Cellule 7 : Définition de l'Augmentation de Données}
Construction d'un sous-modèle séquentiel d'augmentation de données intégrant des transformations géométriques stochastiques (retournement horizontal, rotation aléatoire de 15\%, zoom de 15\% et translation de 10\%) pour prévenir le surapprentissage.

\begin{lstlisting}[language=Python]
data_augmentation = tf.keras.Sequential([
    tf.keras.layers.RandomFlip("horizontal"),
    tf.keras.layers.RandomRotation(0.15),
    tf.keras.layers.RandomZoom(0.15),
    tf.keras.layers.RandomTranslation(
        height_factor=0.1,
        width_factor=0.1
    ),
], name="data_augmentation")
\end{lstlisting}

\begin{figure}[H]
  \centering
  \includegraphics[width=0.82\textwidth]{8.png}
\end{figure}

\subsection{Cellule 8 : Architecture CNN de Référence (Baseline)}
Conception d'un premier modèle convolutif multicouche personnalisé comportant des blocs successifs \texttt{Conv2D} + \texttt{MaxPooling2D} avec des filtres croissants (32, 64, 128), une régularisation par \texttt{Dropout}, un \texttt{GlobalAveragePooling2D} et une couche dense de sortie à 5 unités avec activation softmax.

\begin{lstlisting}[language=Python]
from tensorflow.keras import layers, models

cnn_ameliore = models.Sequential([
    layers.Input(shape=(180, 180, 3)),

    # Diversification des images
    data_augmentation,

    # Extraction des caractéristiques
    layers.Conv2D(32, (3, 3), activation="relu", padding="same"),
    layers.MaxPooling2D((2, 2)),

    layers.Conv2D(64, (3, 3), activation="relu", padding="same"),
    layers.MaxPooling2D((2, 2)),
    layers.Dropout(0.20),

    layers.Conv2D(128, (3, 3), activation="relu", padding="same"),
    layers.MaxPooling2D((2, 2)),
    layers.Dropout(0.30),

    # Classification
    layers.GlobalAveragePooling2D(),
    layers.Dense(128, activation="relu"),
    layers.Dropout(0.40),
    layers.Dense(5, activation="softmax")
])

cnn_ameliore.compile(
    optimizer=tf.keras.optimizers.Adam(learning_rate=0.001),
    loss="sparse_categorical_crossentropy",
    metrics=["accuracy"]
)

cnn_ameliore.summary()
\end{lstlisting}

\begin{figure}[H]
  \centering
  \includegraphics[width=0.82\textwidth]{9.png}
\end{figure}

\subsection{Cellule 9 : Initialisation du Modèle Pré-entraîné MobileNetV2}
Instanciation de l'architecture \textbf{MobileNetV2} entraînée sur le jeu de données ImageNet. La couche sommitale de classification d'origine est exclue (\texttt{include\_top=False}) et les poids pré-entraînés sont gelés (\texttt{trainable=False}) afin de conserver les caractéristiques génériques apprises.

\begin{lstlisting}[language=Python]
from tensorflow.keras.applications import MobileNetV2
from tensorflow.keras.applications.mobilenet_v2 import preprocess_input

# Modèle pré-entraîné, sans sa couche de classification d'origine
base_model = MobileNetV2(
    weights="imagenet",
    include_top=False,
    input_shape=(180, 180, 3)
)

# Geler les poids pré-entraînés au début
base_model.trainable = False
\end{lstlisting}

\begin{figure}[H]
  \centering
  \includegraphics[width=0.82\textwidth]{10.png}
\end{figure}

\subsection{Cellule 10 : Construction du Modèle de Transfer Learning}
Connexion de l'ensemble de la chaîne de traitement : couche d'entrée, bloc d'augmentation, passage à l'échelle des pixels vers $[0, 255]$ pour l'opérateur \texttt{preprocess\_input} de MobileNetV2, passage dans le modèle de base gelé, puis agrégation spatiale (\texttt{GlobalAveragePooling2D}) et tête de classification dense avec \texttt{Dropout} (0.30) et sortie softmax pour les 5 catégories.

\begin{lstlisting}[language=Python]
inputs = layers.Input(shape=(180, 180, 3))

x = data_augmentation(inputs)

# Passage de [0, 1] à [0, 255]
x = layers.Rescaling(255.0)(x)

# Prétraitement spécifique à MobileNetV2
x = layers.Lambda(preprocess_input)(x)

# Extraction des caractéristiques pré-apprises
x = base_model(x, training=False)

# Classification des cinq fleurs
x = layers.GlobalAveragePooling2D()(x)
x = layers.Dropout(0.30)(x)
x = layers.Dense(128, activation="relu")(x)
x = layers.Dropout(0.30)(x)
outputs = layers.Dense(5, activation="softmax")(x)

model_transfer = models.Model(inputs, outputs)

model_transfer.compile(
    optimizer=tf.keras.optimizers.Adam(learning_rate=0.001),
    loss="sparse_categorical_crossentropy",
    metrics=["accuracy"]
)

model_transfer.summary()
\end{lstlisting}

\begin{figure}[H]
  \centering
  \includegraphics[width=0.82\textwidth]{11.png}
\end{figure}

\subsection{Cellule 11 : Entraînement avec Arrêt Précoce (Early Stopping)}
Entraînement supervisé sur 15 époques maximales couplé au callback \texttt{EarlyStopping} surveillant \texttt{val\_loss} avec une patience de 4 époques et restauration automatique des meilleurs poids. Le modèle atteint rapidement une précision de validation supérieure à $92\%$.

\begin{lstlisting}[language=Python]
early_stopping_tl = tf.keras.callbacks.EarlyStopping(
    monitor="val_loss",
    patience=4,
    restore_best_weights=True
)

history_transfer = model_transfer.fit(
    train_ds,
    validation_data=val_ds,
    epochs=15,
    callbacks=[early_stopping_tl]
)
\end{lstlisting}

\begin{figure}[H]
  \centering
  \includegraphics[width=0.82\textwidth]{12.png}
\end{figure}

\subsection{Cellule 12 : Sauvegarde du Modèle Entraîné}
Sérialisation du modèle au format natif moderne \texttt{.keras}, préservant à la fois l'architecture du réseau, les poids optimisés et la configuration de compilation.

\begin{lstlisting}[language=Python]
model_transfer.save("model_transfer.keras")
\end{lstlisting}

\begin{figure}[H]
  \centering
  \includegraphics[width=0.82\textwidth]{13.png}
\end{figure}

\subsection{Cellule 13 : Téléchargement du Modèle Sérialisé}
Utilisation du module \texttt{google.colab.files} pour automatiser le téléchargement direct du fichier binaire du modèle vers le poste client.

\begin{lstlisting}[language=Python]
from google.colab import files

files.download("model_transfer.keras")
\end{lstlisting}

\begin{figure}[H]
  \centering
  \includegraphics[width=0.82\textwidth]{14.png}
\end{figure}

\newpage
\section{Déploiement de l'Application Interactive (\texttt{app.py})}
Pour rendre le modèle accessible aux utilisateurs finaux, une application web réactive et intuitive a été développée à l'aide de la bibliothèque \textbf{Streamlit}.

\subsection{Architecture et Fonctionnalités de l'Application}
L'application \texttt{app.py} intègre l'ensemble du cycle d'inférence en production :
\begin{itemize}
    \item \textbf{Reconstruction et mise en cache du modèle :} utilisation de la fonction \texttt{build\_model()} et du décorateur \texttt{@st.cache\_resource} pour instancier l'architecture exacte et charger les poids (\texttt{model\_transfer.weights.h5}) une seule fois en mémoire vive.
    \item \textbf{Pipeline de prétraitement unifié :} fonction \texttt{preprocess\_image()} assurant la conversion en format RVB, le redimensionnement strict en $180 \times 180$, et la normalisation flottante dans $[0, 1]$.
    \item \textbf{Double mode d'acquisition :} support natif de l'importation de fichiers images (\texttt{jpg}, \texttt{jpeg}, \texttt{png}) ainsi que de la prise de vue en direct par webcam via le composant \texttt{st.camera\_input}.
    \item \textbf{Gestion du seuil de confiance :} curseur dynamique permettant à l'utilisateur de définir le niveau d'incertitude acceptable, affichant une alerte si la confiance maximale est inférieure au seuil choisi.
    \item \textbf{Restitution visuelle des prédictions :} affichage de la classe prédite avec icône représentative, jauge de confiance et histogramme dynamique des probabilités pour chacune des 5 espèces de fleurs.
\end{itemize}

\subsection{Code Source Complet de l'Application (\texttt{app.py})}
Voici l'intégralité du script source développé pour l'interface :

\begin{lstlisting}[language=Python]
import numpy as np
import pandas as pd
import streamlit as st
import tensorflow as tf
from PIL import Image
from tensorflow.keras import layers, models
from tensorflow.keras.applications import MobileNetV2
from tensorflow.keras.applications.mobilenet_v2 import preprocess_input

# ------------------------------------------------------------------
# Configuration
# ------------------------------------------------------------------
IMG_SIZE = (180, 180)
# Meme ordre que metadata.features["label"].names dans TFDS
CLASS_NAMES = ["dandelion", "daisy", "tulips", "sunflowers", "roses"]
CLASS_NAMES_FR = {
    "dandelion": "Pissenlit(dandelion) 🌼",
    "daisy": "Marguerite(daisy) 🌼",
    "tulips": "Tulipe(tulips) 🌷",
    "sunflowers": "Tournesol(sunflowers) 🌻",
    "roses": "Rose(roses) 🌹",
}

st.set_page_config(page_title="Classification de fleurs", page_icon="🌸", layout="centered")


# ------------------------------------------------------------------
# Chargement du modele (mis en cache pour ne le charger qu'une fois)
# ------------------------------------------------------------------
def build_model():
    """Reconstruit exactement l'architecture du notebook."""
    data_augmentation = tf.keras.Sequential([
        layers.RandomFlip("horizontal"),
        layers.RandomRotation(0.15),
        layers.RandomZoom(0.15),
        layers.RandomTranslation(height_factor=0.1, width_factor=0.1),
    ], name="data_augmentation")

    base_model = MobileNetV2(weights=None, include_top=False, input_shape=(180, 180, 3))
    base_model.trainable = False

    inputs = layers.Input(shape=(180, 180, 3))
    x = data_augmentation(inputs)
    x = layers.Rescaling(255.0)(x)
    x = layers.Lambda(preprocess_input)(x)
    x = base_model(x, training=False)
    x = layers.GlobalAveragePooling2D()(x)
    x = layers.Dropout(0.30)(x)
    x = layers.Dense(128, activation="relu")(x)
    x = layers.Dropout(0.30)(x)
    outputs = layers.Dense(5, activation="softmax")(x)
    return models.Model(inputs, outputs)


@st.cache_resource
def load_model():
    model = build_model()
    model.load_weights("model_transfer.weights.h5")
    return model


model = load_model()


# ------------------------------------------------------------------
# Pretraitement (identique au notebook)
# ------------------------------------------------------------------
def preprocess_image(pil_image: Image.Image) -> np.ndarray:
    img = pil_image.convert("RGB")
    arr = np.array(img)
    arr = tf.image.resize(arr, IMG_SIZE)          # redimensionnement 180x180
    arr = tf.cast(arr, tf.float32) / 255.0        # [0,255] -> [0,1]
    return tf.expand_dims(arr, axis=0).numpy()    # (1, 180, 180, 3)


def predict(pil_image: Image.Image):
    batch = preprocess_image(pil_image)
    probs = model.predict(batch, verbose=0)[0]
    return probs


# ------------------------------------------------------------------
# Interface
# ------------------------------------------------------------------
st.title("🌸 Classification de fleurs")
st.write(
    "Charge une photo (ou prends-en une) et le modele **MobileNetV2** "
    "(transfer learning) devine l'espece parmi 5 classes."
)

with st.sidebar:
    st.header("A propos")
    st.write("Dataset : **TF Flowers** (3670 images)")
    st.write("Classes : " + ", ".join(CLASS_NAMES_FR.values()))
    threshold = st.slider("Seuil de confiance", 0.0, 1.0, 0.5, 0.05)

source = st.radio("Source de l'image", ["📁 Fichier", "📷 Camera"], horizontal=True)

if source == "📁 Fichier":
    uploaded = st.file_uploader("Choisis une image", type=["jpg", "jpeg", "png"])
else:
    uploaded = st.camera_input("Prends une photo")

if uploaded is not None:
    image = Image.open(uploaded)
    col1, col2 = st.columns(2)

    with col1:
        st.image(image, caption="Image analysee", use_container_width=True)

    with st.spinner("Analyse en cours..."):
        probs = predict(image)

    best_idx = int(np.argmax(probs))
    best_class = CLASS_NAMES[best_idx]
    confidence = float(probs[best_idx])

    with col2:
        if confidence >= threshold:
            st.success(f"**{CLASS_NAMES_FR[best_class]}**")
        else:
            st.warning(f"Incertain : **{CLASS_NAMES_FR[best_class]}**")
        st.metric("Confiance", f"{confidence * 100:.1f} %")

    st.subheader("Probabilites par classe")
    df = pd.DataFrame({
        "Classe": [CLASS_NAMES_FR[c] for c in CLASS_NAMES],
        "Probabilite": probs,
    }).set_index("Classe")
    st.bar_chart(df)
else:
    st.info("👆 Ajoute une image pour lancer la prediction.")
\end{lstlisting}

\subsection{Images d'Exécution et Démonstration de l'Application Streamlit}
Les figures suivantes illustrent le fonctionnement réel de l'application web Streamlit lors de différentes interactions utilisateur (accueil, téléversement d'image, prédictions d'espèces et capture par webcam).

\subsubsection{Vue d'accueil de l'interface utilisateur}
\begin{figure}[H]
  \centering
  \includegraphics[width=0.82\textwidth]{15.png}
\end{figure}

\subsubsection{Test de classification avec importation d'image (Marguerite / Daisy)}
\begin{figure}[H]
  \centering
  \includegraphics[width=0.82\textwidth]{16.png}
\end{figure}

\subsubsection{Test de classification avec une seconde espèce (Rose / Tulipe)}
\begin{figure}[H]
  \centering
  \includegraphics[width=0.82\textwidth]{17.png}
\end{figure}

\subsubsection{Utilisation de la prise de vue par caméra et gestion du seuil}
\begin{figure}[H]
  \centering
  \includegraphics[width=0.82\textwidth]{18.png}
\end{figure}

\newpage
\section{Conclusion et Perspectives}
Ce projet a permis de concrétiser une chaîne d'ingénierie complète en vision par ordinateur :
\begin{itemize}
    \item L'utilisation du \textbf{Transfer Learning} avec le backbone \textbf{MobileNetV2} a démontré une supériorité nette par rapport à un CNN entraîné de zéro (\textit{scratch}), tant en vitesse de convergence qu'en capacité de généralisation (précision supérieure à 92\% sur l'ensemble de validation).
    \item La régularisation via l'augmentation de données (\texttt{RandomFlip}, \texttt{RandomRotation}, etc.) et le mécanisme de rappel \texttt{EarlyStopping} ont efficacement évité le surapprentissage.
    \item Le déploiement sous \textbf{Streamlit} offre une expérience utilisateur moderne, fluide et réactive, permettant une interaction directe avec le modèle prédictif via des photos locales ou une caméra connectée.
\end{itemize}

\end{document}

